\documentclass{report}
\usepackage{amsmath,amssymb}
\setlength{\parindent}{0mm}
\setlength{\parskip}{1em}
\begin{document}
\begin{center}
\rule{6in}{1pt} \
{\large Misha Kilmer \\
{\bf Tomographic Image Reconstruction via Tensor-Based Dictionary Learning}}

503 Boston Ave  \\ Bromfield-Pearson Bldg \\ Tufts University \\ Medford \\ MA 02155
\\
{\tt misha.kilmer@tufts.edu}\\
Sara Soltani\\
Per Christian Hansen\end{center}

In this talk, we will consider reconstruction of Xray CT images using
priors in the form of a dictionary learned from training images. The
reconstruction process has two stages: a) construction of a tensor
dictionary prior from our training data; b) solution of the
reconstruction problem by optimizing for the expansion tube fiber
coefficients in that dictionary. Our approach differs from past
approaches in that i) we use a third-order tensor representation for our
images and ii) we recast the reconstruction problem using the tensor
formulation. The dictionary learning problem is presented as a
non-negative tensor factorization problem with sparsity constraints. The
reconstruction problem is formulated in a convex optimization framework
by looking for a solution with a sparse representation in the tensor
dictionary. Numerical results show that our tensor formulation leads to
very sparse representations of both the training images and the
reconstructions due to the ability of representing repeated features
compactly in the dictionary.


\end{document}
